\bigskip
\noindent Autres pistes de l'onglet Findings pour ce thème, dans des cotes qui n'ont pas de section ici :\par

\begin{encadre}
\begin{theoreme*}[Grothendieck, pages 93 et 94, amplitude par la caractéristique d'Euler\piste{14-nakai-euler-characteristic-criterion}]\nonverifie
Un faisceau inversible $L$ sur un schéma projectif est ample dès que $\chi(\mathcal O_Z \otimes L^{\otimes n}) \to +\infty$ pour tout sous-schéma fermé intègre $Z$ de dimension $\geq 1$, et de même avec $\dim H^0$ à la place de $\chi$. La page attribue l'énoncé à Mumford et Nakai.
\end{theoreme*}
\end{encadre}
{\small Les feuillets en esquissent une démonstration incomplète (page 94). Ses deux lemmes y sont renvoyés à une \q{lettre de Mumford}. Mathlib n'a pas la caractéristique d'Euler des faisceaux cohérents ni l'amplitude. On se reportera aux notes de Grothendieck (\lot{14}{5}{93}).\par}

\begin{encadre}
\begin{theoreme*}[D'après les pages 1 à 22, forme corrigée par l'édition, section canonique par le déterminant\piste{59-determinant-canonical-section}]\nonverifie
Soient $X/S$ propre et plat, avec $\mathrm{Pic}^0_{X/S} = B$ schéma abélien de dual $A$, et $\delta$ une section de $\mathrm{NS}_{X/S}$. Le déterminant de la cohomologie du faisceau de Poincaré donne une trivialisation canonique du $A$-torseur $\mathrm{Alb}^{\chi(\delta)}_{X/S} \times^A \mathrm{Pic}^{\varphi(\delta)}_{B/S} \times^A (\widetilde{\varphi(\delta)}_*\mathrm{Pic}^\delta_{X/S})^{(-1)}$, où $\varphi : \mathrm{NS}_{X/S} \to \mathrm{NS}_{B/S}$ est une application polynomiale de degré $n - 1$, non additive.
\end{theoreme*}
\end{encadre}
{\small Les feuillets en esquissent une démonstration incomplète (pages 1 à 4 et 16 à 19). La descente y est marquée \q{détailler}, la propriété universelle de $\mathrm{Alb}^1$ annoncée page~22 sans preuve, et l'inverse sur le torseur poussé est de l'édition. Mathlib n'a pas le déterminant de la cohomologie ni le faisceau de Poincaré. On se reportera aux notes de Grothendieck (\lot{59}{1}{1}).\par}

\begin{encadre}
\begin{theoreme*}[Grothendieck, pages 2, 3 et 20, pleine fidélité et composantes connexes\piste{100-full-faithfulness-connected-components}]\nonverifie
Soient $S$ un schéma localement noethérien et $U \subset S$ un ouvert tel que toutes les caractéristiques résiduelles de $S - U$ soient égales à $p$. Le foncteur $A \mapsto (A_U, A[p^\infty], \mathrm{id})$ des schémas abéliens sur $S$ vers les triplets est pleinement fidèle dès que $U$ rencontre chaque composante connexe de $S$, sans hypothèse de normalité, de réduction ni de densité schématique.
\end{theoreme*}
\end{encadre}
{\small Les feuillets en esquissent une démonstration incomplète (page 3). Le lemme d'existence de $S_v$ y renvoie à \q{Invariants 1.2}. Mathlib n'a pas les schémas abéliens ni les groupes $p$-divisibles. On se reportera aux notes de Grothendieck (\lot{100}{1}{2}).\par}

\begin{encadre}
\begin{theoreme*}[Grothendieck, pages 22 à 24, descente des morphismes le long d'un épimorphisme propre\piste{100-proper-epimorphism-hom-descent}]\nonverifie
Soit $f : S' \to S$ un morphisme propre de schémas localement noethériens tel que $\mathcal O_S \to f_*\mathcal O_{S'}$ soit injectif, sans hypothèse de platitude, et soit $S'' = S' \times_S S'$. Pour tout $S$-schéma $X$ non ramifié et localement de type fini, la suite $\mathrm{Hom}_S(S, X) \to \mathrm{Hom}_S(S', X) \rightrightarrows \mathrm{Hom}_S(S'', X)$ est exacte ; le changement de base le long de $f$ est donc pleinement fidèle sur les schémas abéliens.
\end{theoreme*}
\end{encadre}
{\small Démonstration donnée par l'édition, non vérifiée en Lean. Mathlib n'a pas la descente des morphismes le long d'un morphisme propre. On se reportera aux notes de Grothendieck (\lot{100}{2}{22}).\par}

\begin{encadre}
\begin{theoreme*}[Grothendieck, pages 30 à 38, algébrisation par les courbes\piste{100-formal-abelian-curve-criterion}]\nonverifie
Soit $R$ un anneau local noethérien complet. Un schéma abélien formel sur $\operatorname{Spf} R$ est algébrisable si et seulement si son changement de base à $\operatorname{Spf} B'$ l'est pour tout quotient intègre $B$ de dimension un de $R$, de normalisé $B'$, ou encore s'il le devient après un morphisme propre surjectif $S' \to \operatorname{Spec} R$.
\end{theoreme*}
\end{encadre}
{\small Les feuillets en esquissent une démonstration incomplète (pages 32 à 38). Mathlib n'a pas l'algébrisation des schémas formels. On se reportera aux notes de Grothendieck (\lot{100}{2}{30}).\par}

\begin{encadre}
\begin{theoreme*}[Grothendieck, pages 39 à 48, faisceau des quasi-fonctions\piste{130-quasi-function-sheaf}]\nonverifie
Les fonctions de Weil forment un faisceau de Zariski $QF_X$, somme amalgamée du faisceau $\Phi_X$ des fonctions admissibles et de $\mathcal R^*_X$ au-dessus de $\mathcal O^*_X$, dans une suite exacte $0 \to \Phi_X \to QF_X \to \mathrm{Div}_X \to 0$. Pour $X$ quasi-projectif, tout diviseur de Cartier est le diviseur d'une quasi-fonction globale.
\end{theoreme*}
\end{encadre}
{\small Les feuillets en esquissent une démonstration incomplète (pages 41 à 48). La page s'arrête avant le cas de $\mathbf P^r$. Mathlib n'a pas les fonctions de Weil ni les diviseurs de Cartier. On se reportera aux notes de Grothendieck (\lot{130}{2}{39}).\par}

\begin{encadre}
\begin{theoreme*}[Grothendieck, pages 22 à 28, déterminants des parties étales\piste{9-det-etale-tate-iso}]\nonverifie
Pour un schéma abélien $A$ ordinaire localement polarisable sur une base connexe de caractéristique~$p$, $\alpha_A = \det T_p(\tilde\varphi)_{\text{ét}}/\sqrt{\deg\tilde\varphi}$ ne dépend pas de la polarisation $\varphi$ et est un isomorphisme de réseaux $\det T_p(A)_{\text{ét}} \simeq \det T_p(A^*)_{\text{ét}}$, compatible au changement de base ; pour une isogénie $u$, on a seulement $\det T_p(u)_{\text{ét}} \cdot \det T_p(u^*)_{\text{ét}} = \deg u$.
\end{theoreme*}
\end{encadre}
{\small Les feuillets en esquissent une démonstration incomplète : définitions, énoncés a) et b) et leurs arguments sont sur les pages, mais l'identité $\Delta(u)\Delta(u^*) = \deg u$ y est introduite par une justification illisible (la page~25, qui l'attribuait à un théorème $p$-adique, est annulée), les points c) et d) sont laissés \q{à vérifier}, et les hypothèses de polarisabilité locale et de base connexe sont de l'édition.\par}

\begin{encadre}
\begin{theoreme*}[Grothendieck, pages 24 et 26, déterminant et racine du degré\piste{9-symmetric-positive-det-sqrt-degree}]\nonverifie
Pour une variété abélienne ordinaire polarisée $A$ et $u \in \mathrm{End}(A) \otimes \mathbf Q$ symétrique et positif pour l'involution de Rosati, le déterminant de $u$ sur la partie étale du module de Tate $p$-adique vaut $+\sqrt{\deg u}$.
\end{theoreme*}
\end{encadre}
{\small Les feuillets en esquissent une démonstration incomplète : la page~24 ramène l'énoncé au signe, et la page~26 l'obtient par positivité, $u = v^2$ sur le cône positif ; l'identité $\Delta(u)\Delta(u^*) = \deg u$ y repose sur une justification illisible, et la restriction de l'argument de rationalité aux éléments symétriques est de l'édition.\par}

\begin{encadre}
\begin{theoreme*}[Grothendieck, pages 23 à 29 et 33, équivalence d'Albanese\piste{14-albanese-equivalence-ring}]\nonverifie
Sur les cycles de toute codimension, posons $z \sim 0$ quand $z = Z_*(t)$ pour un cycle $Z$ sur $X \times T$, $T$ projective lisse connexe, et un $0$-cycle $t$ de degré~$0$ sur $T$ tel que $\mathrm{alb}_T(t) = 0$. Cette relation est comprise entre les équivalences rationnelle et algébrique, passe aux produits, aux images réciproques et aux images directes propres, coïncide avec l'équivalence rationnelle sur les diviseurs et avec le degré et la somme d'Albanese sur les $0$-cycles ; dans l'anneau quotient, le produit de deux classes algébriquement triviales est nul.
\end{theoreme*}
\end{encadre}
{\small Les feuillets en esquissent une démonstration incomplète : les théorèmes~1.4 et~1.6 sont démontrés (pages 28 et~29), mais la proposition~1.2~(iii) est laissée \q{à vérifier} (page~27), le~(iv) est fait \q{par l'AQT}, abréviation non reconstituée, et le corollaire~2.4 (page~33) est sans démonstration.\par}

\begin{encadre}
\begin{theoreme*}[Grothendieck, pages 9 et 10, Riemann--Roch et transformation de Fourier sur une base\piste{47-chow-fourier-todd-over-base}]\nonverifie
Sur une base $S$, Grothendieck--Riemann--Roch pour $m_A$ et $\mathrm{pr}_2$ donne $\mathrm{ch}(x * y) = (\mathrm{ch}\,x * \mathrm{ch}\,y)\cdot\mathrm{Todd}(t_A)$ sur un schéma abélien $A$ et, pour $\Psi_D = \mathrm{pr}_{2*}(\mathrm{pr}_1^*(-)\cdot\exp D)$, $\Psi_D(a * b) = \Psi_D(a)\,\Psi_D(b)$ sans correction ; la formule qui échange cup-produit et produit de Pontrjagin prend un facteur venu de~$S$.
\end{theoreme*}
\end{encadre}
{\small La page~9 affirme $(\ast)$ par Riemann--Roch (\q{On vérifie}) sans le calcul, et la page~10 dérive la seconde formule avec le seul facteur $\mathrm{Todd}(t_A)$ ; l'édition corrige ce facteur en $\mathrm{Todd}(t_A)\,\mathrm{Todd}(t_B)\,\mathrm{ch}(\omega_A)$, la seconde application de $(\ast)$ se faisant sur $B$. Le facteur venu de $S$ de l'encadré est cette forme corrigée, qui est de l'édition. Les pages~3 à~7 qui précèdent établissent l'inversion de Fourier relative, $\Phi_{{}^sL} \circ \Phi_L \simeq s_A^*(-) \otimes \tau_A[-n]$ sur les complexes parfaits, $\tau_A$ le déterminant de l'algèbre de Lie de $A$ : c'est le théorème~1.1 de Mukai (1987) \cite{mukai1987}, pour un schéma abélien sur une base quelconque, avec la torsion par le dual du fibré canonique relatif. Deninger et Murre (1991) \cite{deningermurre1991} étudient la transformation de Fourier sur l'anneau de Chow rationnel d'un schéma abélien ; la comparaison de leurs formules avec le facteur corrigé n'a pas été faite. On applique ici à $m_A$ et à $\mathrm{pr}_2$ le théorème de Riemann--Roch de Grothendieck, qui n'est vérifié en Lean nulle part dans l'annexe ; la partie \q{K-théorie, classes de Chern et Riemann--Roch} (section~\ref{sec:2}) n'en contrôle que quelques calculs.\par}

\begin{encadre}
\begin{theoreme*}[Grothendieck, pages 16, 17 et 20, bonne réduction de l'Albanese\piste{51-albanese-good-reduction-codim-two}]\nonverifie
Si $X$ est projectif et plat sur un anneau de valuation discrète $V$, à fibres lisses hors d'un fermé de codimension $\geq 2$, de fibre générique géométriquement intègre et normale, $\mathrm{Alb}(X_K)$ et $(\mathrm{Pic}^0_{X_K})_{\mathrm{red}}$ ont bonne réduction.
\end{theoreme*}
\end{encadre}
{\small L'hypothèse que la fibre générique est normale est de l'édition. Les feuillets en esquissent une démonstration incomplète : une courbe $C$ lisse sur $V$ qui engendre la variété abélienne, puis Koizumi--Shimura ; l'existence de $C$ et le fait qu'elle engendre sont affirmés (\q{on sait}) sans démonstration.\par}

\begin{encadre}
\begin{theoreme*}[Grothendieck, pages 98 et 99, section hyperplane générique\piste{57-generic-hyperplane-h1-injective}]\nonverifie
Soient $X \subset \mathbf P^r$ normal intègre de dimension $\geq 2$ sur un corps algébriquement clos et $Y_{\bar K}$ sa section hyperplane générique sur la clôture algébrique de $k(t_1, \ldots, t_r)$. Alors $H^1(X, G) \to H^1(Y_{\bar K}, G)$ est injectif pour $G = \mathbf G_a$ et pour $G$ fini commutatif, $\mu_p$ et $\alpha_p$ compris ; le noyau de $\mathrm{Pic}^\tau_{X_{\bar K}} \to \mathrm{Pic}^\tau_{Y_{\bar K}}$ est un $p$-groupe unipotent fini.
\end{theoreme*}
\end{encadre}
{\small Les feuillets en esquissent une démonstration incomplète : la page~98 ramène le théorème aux corollaires 1 à~4, et la page~99 plonge $H^1(X, \alpha_p)$ et $\mathrm{Pic}(X)[p]$ dans $H^0(X, \tilde\Omega^1)$ par Cartier ; l'injectivité de la restriction sur $H^0(\tilde\Omega^1)$, dite \q{presque triviale}, n'est pas faite, et le corollaire~1 n'est pas achevé.\par}

\begin{encadre}
\begin{theoreme*}[Grothendieck, pages 226 à 228 et 247 à 271, Picard analytique et complété\piste{57-analytic-pic-completion-criterion}]\nonverifie
Soient $f : X \to Y$ un morphisme propre d'espaces analytiques complexes et $y \in Y$. L'application $(R^1f_*\mathcal O^*_X)_y \to \varprojlim \mathrm{Pic}(X_n)$ est un isomorphisme si et seulement si $(R^1f_*\mathcal O_X)_y$ est de longueur finie ; le système $(\mathrm{Pic}(X_n))$ est de Mittag-Leffler ; $X$ est projectif au voisinage de $y$ si et seulement si tous les $X_n$ sont projectifs.
\end{theoreme*}
\end{encadre}
\begin{demopage}[pages 226 à 228 et 247 à 261]
La lettre à Hironaka du 6~juillet 1962 (pp.~226--228) énonce le théorème et en annonce la démonstration par les théorèmes de Grauert et la suite exponentielle ; le brouillon la donne. Par Grauert (lemme~1.3), $(H^i(X_n, \mathcal O_{X_n}))_n$ est de Mittag-Leffler, de limite $(R^if_*\mathcal O_X)^\wedge_y$ ; la suite exponentielle, où $H^i(X_n, \mathbf Z) = H^i(X_0, \mathbf Z)$ est de type fini, en tire la même condition pour $(H^i(X_n, \mathcal O^*_{X_n}))_n$ et une suite exacte (proposition~1.5). Le lemme des cinq, qui compare celle-ci à la suite des tiges, rend la flèche de $(R^if_*\mathcal O^*_X)_y$ vers la limite injective, et bijective si et seulement si $(R^if_*\mathcal O_X)_y$ est de longueur finie : c'est le théorème pour $i = 1$ (remarque~1.7, et corollaire~1.8 pour la projectivité).
\end{demopage}

\begin{encadre}
\begin{exemple*}[Grothendieck, pages 11 à 13, torseur de Picard en dimension 2\piste{59-surface-even-chi-torsor}]\nonverifie
Pour un schéma abélien $A/S$ de dimension relative~$2$ et une section $\delta$ de $\mathrm{NS}_{A/S}$ avec $\chi(\delta)$ pair, le torseur $\mathrm{Pic}^\delta_{A/S}$ sous le dual est trivial.
\end{exemple*}
\end{encadre}
{\small La page le donne en exemple d'un calcul (\q{Exemple}, page~13). Les feuillets en esquissent une démonstration incomplète : la page~13 déduit la trivialité des puissances $\chi^{n-2}(\chi - 1)$ et $2\chi^{n-2}$ du torseur, et conclut pour $n = 2$ sur un \q{si} de lecture incertaine ; les compatibilités des trivialisations que les marges demandent ne sont pas démontrées.\par}

\begin{encadre}
\begin{theoreme*}[Pages 36 à 38, forme corrigée par l'édition, points de dimension un\piste{100-baire-dimension-one-points}]\nonverifie
Dans le spectre d'un anneau local noethérien complet, une famille dénombrable de fermés dont la réunion contient tous les points $x$ tels que $\dim \overline{\{x\}} \leq 1$ a une sous-famille finie qui recouvre le spectre.
\end{theoreme*}
\end{encadre}
{\small La page barre \q{complet}, que l'édition rétablit : sans cette hypothèse l'énoncé est faux. Les feuillets démontrent le cas de dimension $\leq 1$ et le corollaire, une propriété de Baire sur les $\mathfrak m^n$ ; le cas général de la proposition s'interrompt page~36 sur des lignes illisibles et annulées, et le passage du corollaire à la proposition n'est pas écrit.\par}

\begin{encadre}
\begin{theoreme*}[Grothendieck, pages 47 à 52, prolongement vers un modèle de Néron\piste{100-neron-extension-p-divisible}]\nonverifie
Sur un trait $S$, pour tout modèle de Néron $G$ et tout groupe $p$-divisible $M$ sur $S$, tout homomorphisme $M_K \to G_K$ des fibres génériques se prolonge de façon unique en $M \to G$, le théorème de pleine fidélité de Tate pour les groupes $p$-divisibles sur $S$ étant admis.
\end{theoreme*}
\end{encadre}
{\small La conclusion repose sur le théorème de Tate, et le cas où $G_K$ est extension d'une variété abélienne par un tore est traité dans le tapuscrit des pages 47 à~50, dont l'auteur n'est pas établi ; la part des feuillets manuscrits est le lemme des pages 51 et~52. Les feuillets en esquissent une démonstration incomplète : ce lemme renvoie deux fois à \q{on sait}, et la page~52 l'applique au modèle de Néron et conclut \q{On gagne !}.\par}

\begin{encadre}
\begin{conjecture*}[Grothendieck, pages 7 et 8, intégralité sur une variété abélienne\piste{47-dual-polarisation-divisibility}]
Sur une variété abélienne, pour un fibré en droites $L$ de classe de Néron--Severi $\delta$ avec $\chi = \chi(\delta) \neq 0$ et $\delta'$ la classe déterminée par $\det \Phi_L(L)$, les classes $\binom{\chi}{k}\delta'^k/\chi^m$ sont entières, comme conséquences de $c(\Phi_L(\mathrm{cl}\,L)) = (1 - D'/\chi)^\chi$.
\end{conjecture*}
\end{encadre}
{\small La lecture prend $m = k$, que la ligne du caractère de Chern de la page impose, et en donne une démonstration sur un corps algébriquement clos par l'image réciproque le long de l'isogénie $\varphi_L$, injective sur l'anneau de Chow rationnel ; l'exposant $m = k - 2$ de la page donne un énoncé plus faible, vrai si celui de la lecture l'est. Cette démonstration est de l'édition et n'est pas vérifiée en Lean ; l'argument surchargé de $\varphi_L$ reste à relire sur le fac-similé. Mathlib n'a pas les classes de Chern sur les variétés abéliennes. On se reportera aux notes de Grothendieck (\lot{47}{1}{7}).\par}

\begin{encadre}
\begin{theoreme*}[D'après la page 3, forme reconstituée par l'édition, critère de platitude\piste{49-flatness-lemma-pushforward}]\nonverifie
Soient $S$ localement noethérien, $\varphi : X \to Y$ un morphisme de $S$-schémas de type fini, $F$ quasi-cohérent sur $X$ et plat sur $S$, $G$ cohérent sur $Y$ et $G \to \varphi_*F$. Si $G \to \varphi_*F$ est injectif sur la fibre en $s$ aux points au-dessus de $y$, $G$ est plat sur $S$ en $y$.
\end{theoreme*}
\end{encadre}
{\small L'édition retient la forme barrée de la clause~c), la clause conservée étant coupée par un mot illisible. Mathlib n'a pas la platitude des faisceaux cohérents sur les schémas. On se reportera aux notes de Grothendieck (\lot{49}{1}{3}).\par}

\begin{encadre}
\begin{theoreme*}[Grothendieck, pages 42 à 49, rigidité des sous-groupes\piste{57-rigid-subgroups-one-fibre}]\nonverifie
Soient $S$ localement noethérien connexe et $G$ un schéma en groupes de type fini sur $S$. Deux sous-schémas en groupes fermés \q{rigides} de $G$, c'est-à-dire commutatifs, plats, et dont chaque fibre géométrique est l'adhérence schématique de sa torsion d'ordre premier à la caractéristique résiduelle, qui coïncident sur une fibre, sont égaux.
\end{theoreme*}
\end{encadre}
{\small Mathlib n'a pas les schémas en groupes sur une base noethérienne. On se reportera aux notes de Grothendieck (\lot{57}{3}{42}).\par}

\begin{encadre}
\begin{theoreme*}[Grothendieck, pages 3 à 5, symbole local de Néron\piste{130-neron-symbol-component-correction}]\nonverifie
Pour une variété abélienne $A$ sur le corps des fractions $K$ d'un anneau de valuation discrète, le symbole local de Néron $(X, a)_v$ est la somme d'un nombre d'intersection $i(\bar X, \bar a)$ sur le modèle de Néron et d'une correction $v(X, a)$ caractérisée par quatre axiomes (bilinéarité, invariance par translation, normalisation par la partie verticale de $\mathrm{div}\,\bar f$, bornitude), avec $m\,v(X, a) \in \mathbf Z$ pour $m$ l'ordre du groupe des composantes ; d'où un accouplement $A^\vee(K) \times A(K) \to \mathbf Z/m\mathbf Z$.
\end{theoreme*}
\end{encadre}
{\small Mathlib n'a pas les modèles de Néron ni les symboles locaux. On se reportera aux notes de Grothendieck (\lot{130}{1}{3}).\par}
